\documentclass[11pt]{article}

\oddsidemargin=17pt \evensidemargin=17pt
\headheight=9pt     \topmargin=26pt
\textheight=564pt   \textwidth=433.8pt

\usepackage{amsmath, amsthm, amssymb,enumitem}
\setlist[enumerate]{parsep=0pt plus 4pt,topsep=3pt plus 4pt}
\setlist[itemize]{parsep=0pt plus 4pt,topsep=3pt plus 4pt,itemsep=.5ex}

\leftmargini=5.5ex
\leftmarginii=3.5ex

%For separated lists with consecutive numbering
\newcounter{separated}
	
\newcommand{\excise}[1]{}
\newcommand{\comment}[1]{{$\star$\sf\textbf{#1}$\star$}}
\newcommand{\exnumber}[1]{\noindent\textbf{#1}}
	
%new math symbols taking no arguments
\newcommand\0{\mathbf{0}}
\newcommand\CC{\mathbb{C}}
\newcommand\FF{\mathbb{F}}
\newcommand\NN{\mathbb{N}}
\newcommand\QQ{\mathbb{Q}}
\newcommand\RR{\mathbb{R}}
\newcommand\ZZ{\mathbb{Z}}
\newcommand\bb{\mathbf{b}}
\newcommand\kk{\Bbbk}
\newcommand\mm{\mathfrak{m}}
\newcommand\xx{\mathbf{x}}
\newcommand\yy{\mathbf{y}}
\newcommand\GL{\mathit{GL}}
\newcommand\pset{\mathcal{P}}
\newcommand\into{\hookrightarrow}
\newcommand\nsub{\trianglelefteq}
\newcommand\onto{\twoheadrightarrow}
\newcommand\minus{\smallsetminus}
\newcommand\goesto{\rightsquigarrow}

% theorem environment
\newtheorem*{theorem}{Theorem}

%redefined math symbols taking no arguments
\newcommand\<{\langle}
\renewcommand\>{\rangle}
\renewcommand\iff{\Leftrightarrow}
\renewcommand\phi{\varphi}
\renewcommand\implies{\Rightarrow}
	
%new math symbols taking arguments
\newcommand\ol[1]{{\overline{#1}}}
	
%redefined math symbols taking arguments
\renewcommand\mod[1]{\ (\mathrm{mod}\ #1)}
	
%roman font math operators
\DeclareMathOperator\im{im}
	
%for easy 2 x 2 matrices
\newcommand\twobytwo[1]{\left[\begin{array}{@{}cc@{}}#1\end{array}\right]}
	
%for easy column vectors of size 2
\newcommand\tworow[1]{\left[\begin{array}{@{}c@{}}#1\end{array}\right]}
	
	
%%%%%%%%%%%%%%%%%%%%%%%%%%%%%%%%%%%%%%%%%%%%%%%%%%%%%%%%%%%%%%%%%%%%%%
\begin{document}%%%%%%%%%%%%%%%%%%%%%%%%%%%%%%%%%%%%%%%%%%%%%%%%%%%%%%
%%%%%%%%%%%%%%%%%%%%%%%%%%%%%%%%%%%%%%%%%%%%%%%%%%%%%%%%%%%%%%%%%%%%%%


\title{\mbox{}\\[-8ex]Relations Solutions\normalsize
	\\[-2.5ex]}
\author{Solutions by: Raindrops\_drop \\[1ex]
	\\[-1ex]}
\date{September 12, 2026} 
\maketitle

\vspace{-3ex}%
\noindent

\exnumber{Exercises}

\exnumber{1.} Define two points $(x_0, y_0)$ and $(x_1, y_1)$ of the plane to be equivalent
if $y_0 - x_0^2 = y_1 - x_1^2$. Check that this is an equivalence relation and describe
the equivalence classes.
\begin{proof}
	For the relation described in this problem to be an equivalence relation, it must
	be reflexive, symmetric, and transitive. Let $(x, y) \in \RR^2$, and note 
	$y - x^2 = y - x^2$ since $x$ and $y$ are the same. Therefore, the relation is
	reflexive. Let $(x_0, y_0), (x_1, y_1) \in \RR^2$ such that $(x_0, y_0) = (x_1, y_1)$, thus
	$y_0 - x_0^2 = y_1 - x_1^2$. Note that "=" as used in this equality for 
	real numbers is symmetric and therefore $y_1 - x_1^2 = y_0 - x_0^2$. Hence,
	$(x_1, y_1) = (x_0, y_0)$ and we can conclude our relation is symmetric. 
	Let $(x_0, y_0), (x_1, y_1), (x_2, y_2) \in \RR^2$ such that 
	$y_0 - x_0^2 = y_1 - x_1^2$ and $y_1 - x_1^2 = y_2 - x_2^2$. Since "=" as used
	here is the usual equality for real numbers, we note that it is transitive.
	As such, $y_0 - x_0^2 = y_2 - x_2^2$ and we conclude $(x_0, y_0) = (x_2, y_2)$.
	Therefore, our relation is transitive and we can conclude that the relation
	described in this problem is an equivalence relation.
	
	We now examine what the equivalence classes are: They are parabolas. The points on
	a parabola will all be equivalent to each other under this equivalence relation. 
	In particular, the parabola must be of the form $f(x) = x^2 + c$ for some
	constant $c \in \RR$. Note these classes cover all points on the plane and are
	disjoint of each other. Further, that points on parabolas of forms other than 
	this one, are already in some parabola $f(x) = x^2 + c$ and that 
	also have a difference that is dependent on $x$ in some way, i.e. $y - x^2 = a^2 + bx$
	in such case each point on the parabola will have a difference that is
	not equal to another point in the parabola except possibly its reflection across the
	$y$-axis if $b = 0$ and $a \neq 1$, but in such case this pair of points is already
	in another equivalence class, mainly $f(x) = x^2 + c$.
\end{proof}

\exnumber{2.} Let $C$ be a relation on a set $A$. If $A_0 \subseteq A$, define the 
\textbf{\textit{restriction}} of $C$ to $A_0$ to be $C \cap (A_0 \times A_0)$.
Show that the restriction of an equivalence relation is an equivalence relation.

\begin{proof}
	Let $C_0 = C \cap (A_0 \times A_0)$. If $A_0 = \emptyset$, then $C_0 = \emptyset$ also.
	Note if $A$ is a non-empty set, then $C_0$ is not an equivalence relation on $A$ since it 
	is not reflexive. If, however, $A$ is empty, then $C_0$ is vacuously an equivalence relation on
	$A$. Thus, we will assume that $A_0$ is non-empty. We begin by showing that $C_0$ must be reflexive.
	Let $x \in A_0$, then $(x, x) \in A_0 \times A_0$. Since $C$ is an equivalence relation on $A$
	we have that $(x, x) \in C$ by definition. Thus, $(x, x) \in C_0$ and we can conclude that 
	$C_0$ is an reflexive on $A_0$. To show that $C_0$ is symmetric, let $(x, y) \in C_0$. 
	Thus, $(x, y) \in C$ and $(x, y) \in A_0 \times A_0$. This implies that $x, y \in A_0$ and thus
	$(y, x) \in A_0 \times A_0$ also. Furthermore, since $C$ is a relation on $A$, $(y, x) \in C$.
	Thus, $(y, x) \in C_0$ if $(x, y) \in C_0$, showing that $C_0$ is symmetric. Finally, to show that
	$C_0$ is transitive, let $(x,y), (y, z) \in C_0$. Note that $x, y, z \in A_0$ 
	by construction. Hence, $(x, z) \in A_0 \times A_0$. Since $C$ is a relation on $A$, it must 
	be that $(x, z) \in C$. Therefore, $(x, z) \in C_0$, showing that $C_0$ is transitive and thus 
	is an equivalence relation on $A_0$.
\end{proof}

\exnumber{3.} Here is a "proof" that every relation $C$ that is both symmetric and 
transitive is also reflexive: "Since $C$ is symmetric, $aCb$ implies $bCa$. Since 
$C$ is transitive, $aCb$ and $bCa$ together imply $aCa$, as desired." Find the flaw in
this argument.

The main flaw with this argument is that reflexivity requires all elements in a set to be included in the 
relation, i.e., $(x, x) \in R$ for all $x \in A$, while symmetry and transitivity do not have this 
requirement. In particular, say you have the set $A = \{a, b, c, d\}$
with the relation $R = \{(a, b), (b, a), (b, c), (c, b), (a, c), (c, a), (a, a), (b, b), (c, c)\}$.
Note that $R$ is symmetric and transitive; however, it is not reflexive since $(d, d) \notin R$.


\exnumber{4.} Let $f: A \to B$ be a surjective function. Let us define a relation on $A$ 
by setting $a_0 \sim a_1$ if 
$$
f(a_0) = f(a_1).
$$ 
\begin{enumerate}[label=(\alph*)]
	\item Show that this is an equivalence relation.
	\begin{proof}
		Note that for $a \in A$, $f(a) = f(a)$, thus $a \sim a$ for all $a \in A$.
		Let $x, y \in A$ such that $f(x) = f(y)$. Since $=$ is assumed to be an equivalence
		relation on $B$, it must be that $f(y) = f(x)$, thus showing that $y \sim x$ if $x \sim y$.
		Let $x, y, z \in A$ such that $x \sim y$ and $y \sim z$. By definition, we have 
		that $f(x) = f(y)$ and $f(y) = f(z)$. Since $=$ is an equivalence relation on $B$ we have
		that $f(x) = f(z)$. Therefore, $x \sim z$, and we conclude that $\sim$ is an equivalence
		relation on $A$.
	\end{proof}
	\item Let $A^*$ be the set of equivalence classes. Show that there is a bijective 
	correspondence of $A^*$ with $B$.
	\begin{proof}
		Let $g: A^* \to B$ be a function such that $g(a^*) = b$ where $a^*$ is an equivalence
		class in $A$ determined by $a$, and $b = f(a)$. We claim that $g$ is bijective. 
		Consider $b \in B$, since $f$ is surjective, there exists some $a \in A$ such
		that $f(a) = b$. Hence, $f(a*) = b$ by construction and we can conclude that 
		$g$ is surjective. Let $g(x^*) = b$ and $g(y^*) = b$. By construction, $f(x) = b$
		and $f(y) = b$. Note this implies that $x \sim y$ and thus both are elements of the 
		same equivalence class, i.e., $x^* = y^*$ since equivalence classes are either disjoint 
		or equal. Therefore, $g$ is injective, and we can conclude that $g$ is bijective. 
	\end{proof}
\end{enumerate}

\pagebreak

\exnumber{5.} Let $S$ and $S'$ be the following subsets of the plane:
$$
S = \{(x, y) \mid y = x + 1 \text{ and } 0 < x < 2\},
$$
$$
S' = \{(x, y) \mid y - x \text{ is an integer}\}.
$$
\begin{enumerate}[label=(\alph*)]
	\item Show that $S'$ is an equivalence relation on the real line and $S' \supseteq S$.
	Describe the equivalence classes of $S'$.
	
	\begin{proof}
		Let $x \in \RR$ and note that $x - x = 0$ which is an integer. Thus, $xS'x$ for all
		$x \in \RR$. Let $xS'y$, this implies that $y - x = n$ for some integer $n$.Note
		we can rewrite this equation as $x - y = -n$ where $-n$ is still an integer, and thus,
		$yS'x$. Let $xS'y$ and $yS'z$, then $y - x = n$ and $z - y = m$ for integers $n$ and $m$.
		Note these equations imply $z - x = m + n$. Since $n + m$ is an integer, we conclude that 
		$xS'z$. Therefore, $S'$ is an equivalence relation on the real line.
		
		Let $(x, y) \in S$, then $y - x = 1$ where $0 < x <2$. Since $y - x$ is equal to an
		integer, they belong to an equivalence class in $S'$. Thus, $(x, y) \in S'$, and we 
		can therefore conclude that $S \subseteq S'$. 
		
		In regard to the equivalence classes of $S'$, note that all the integers are in the 
		same equivalence class as $0$, since their difference with zero is always an integer. 
		Every rational number between zero and one are in their own distinct equivalence class,
		since rational numbers between zero and one are all less than one away from each other.
		This is also true for the irrational numbers that are between zero and one. In particular,
		the numbers between zero and one, including zero also, can determine the equivalence classes
		of this equivalence relation. As such, we could represent the equivalence classes of 
		this equivalence relation as $[0,1)$. 
	\end{proof}
	
	\item Show that given any collection of equivalence relations on a set $A$, their 
	intersection is an equivalence relation on $A$.
	
	\begin{proof}
		Let $R = \bigcap_{i \in I} R_i$ where $R_i$ are equivalence relations on $A$ and 
		$I$ is some indexing set. Let $x \in A$, since every $R_i$ is an equivalence relation
		it must be that $(x, x) \in R_i$ for all $i \in I$. Thus, $(x, x) \in R$ by definition.
		Let $(x, y) \in R$, this implies that $(x, y) \in R_i$ for all $i$. Since each $R_i$
		is an equivalence relation, $(y, x) \in R_i$ for all $i$. Hence, $(y, x) \in R$.
		Let $(x, y), (y, z) \in R$. This implies that $(x, y), (y, z) \in R_i$ for all $i$.
		Since each $R_i$ is an equivalence relation, it must be that $(x, z) \in R_i$ for all $i$.
		Therefore, $(x, z) \in R$ and we can thus conclude that $R$ is an equivalence relation on
		$A$.
	\end{proof}
	
	\item Describe the equivalence relation $T$ on the real line that is the intersection
	of all equivalence relations on the real line that contain $S$. Describe the 
	equivalence classes of $T$.
	
	\begin{proof}
		One way to find $T$ is to start with the set $S$, and construct an equivalence class around
		$S$. As a first step, we will have that $S \subseteq T$. This is not an equivalence class
		since $S$ is not reflexive, transitive, nor symmetric. Thus, we must include the set
		$S_R = \{(x, x) \mid x \in \RR\}$. Note $S \cup S_R \subseteq T$ is now reflexive.
		To make this relation symmetric, we can look back at the definition of $S$ and construct
		a set that is symmetric to $S$. Recall $S = \{(x, y) \mid y = x + 1 \text{ and } 0 < x < 2\}$.
		A set that is symmetric to $S$ can be defined as $S_S = \{(x, y) \mid y = x - 1 \text{ and }
		1 < x < 3\}$. Thus, the set $S \cup S_R \cup S_S \subseteq T$ is now reflexive and symmetric.
		We need only make this set transitive to truly make it an equivalence relation.
		Note there are elements in $S$ that should provide other points to be included through transitivity,
		like $(0.5, 1.5)$ and $(1.5, 2.5)$. These points should imply that $(0.5, 1.5)$ to also be 
		included. Note that this can be accomplished by including the set $S_{T_1} = \{(x, y)
		\mid y = x + 2 \text{ and } 0 < x < 1\}$. Since we included this set, we must include
		a set symmetric to it, $S_{T_2} = \{(x, y) \mid y = x - 2 \text{ and } 2 < x < 3\}$.
		Thus, note that $S \cup S_R \cup S_S \cup S_{T_1} \cup S_{T_2} = T$ is a reflexive, symmetric,
		and transitive subset of $\RR \times \RR$, and is therefore an equivalence relation. Note $T$
		must be the intersection of all equivalence relations that include $S$ since it is the 
		smallest equivalence relation that includes $S$ by construction.
		
		In regard to the equivalence classes of $T$, note that if $x < 0$ or $x > 0$, its equivalence
		class is the set $\{x\}$ since it is only equivalent to itself by construction. If 
		$0 < x < 1$, then its equivalence class is the set $\{x, x+1, x+2\}$ by construction.
		If $1 < x < 1$, then its equivalence class is the set $\{x-1, x, x+1\}$ by construction.
		Finally, if $2 < x < 3$, then its equivalence class is the set $\{x-2, x-1, x\}$ by construction.
	\end{proof}
	
\end{enumerate}

\exnumber{6.} Define a relation on the plane by setting
$$
(x_0, y_0) < (x_1, y_1)
$$
if either $y_0 - x_0^2 < y_1 - x_1^2$, or $y_0 - x_0^2 = y_1 - x_1^2$ and $x_0 < x_1$. Show
that this is an order relation, and describe it geometrically. 

\begin{proof}
	Let $(x, y), (n, m)$ be distinct points on the plane. This implies that either $x \neq n$, 
	$y \neq m$, or both. If $x \neq n$, either $x^2 < n^2$ or $x^2> n^2$ or $x^2 = n^2$.
	Without loss of generality let $x^2 < n^2$. Since $y = m$, we note that $y - x^2 > m - n^2$.
	Hence $(x, y) > (n, m)$. If $y \neq m$, then either $y < m$ or $y > m$. Without loss of generality
	let $y < m$. Since $x = n$, this implies that $y - x^2 < m - n^2$. Hence, $(x, y) < (n, m)$.
	If both $x = n$ and $y = m$ it may be that $y - x^2 = m - n^2$. Note then that $x \neq n$ and
	thus either $x < n$ or $x > n$, in either case we can still compare both points of the plane.
	Therefore, the relation in this exercise is comparable.
	
	Let $(x, y)$ be a point on the plane. Note that $y - x^2 = y - x^2$ and that $x = x$. Hence,
	our relation for this exercise does not hold for this point and itself. In particular, it implies
	that a point is not comparable with itself. Therefore, our relation is non-reflexive. 
	
	Let $(x, y) < (u, v)$ and $(u, v) < (s, t)$. Then, by definition, $y - x^2 < v - u^2$ or 
	$y - x^2 = v - u^2$ and $x < u$. Similarly, $v - u^2 < t - s^2$ or $v - u^2 = t - s^2$ and 
	$u < s$. Without loss of generality let $y - x^2 < v - u^2$ and $v - u^2 < t - s^2$.
	This implies that $y - x^2 < t - s^2$, showing that $(x, y) < (s, t)$ as desired. Therefore,
	this relation is transitive and we can conclude that this is an order relation.
	
	Geometrically, we note that every point on the plane will lie in exactly one parabola of 
	the form $y = x^2 + c$ for some $c \in \RR$. In particular, by our choice of relation we have
	made it so that if one point lies in a parabola higher up than another, then that point will be
	greater than the other. If two points lie on the same parabola, then the point most to the right will
	be the greater of the two.   
\end{proof}

\exnumber{7.} Show that the restriction of an order relation is an order relation.

\begin{proof}
	Let $C$ be an order relation on the set $A$, and $A_0 \subseteq A$ be a subset of $A$. Consider the 
	relation $C_0 = C \cap (A_0 \times A_0)$. If $A_0 = \emptyset$ note that $C_0 = \emptyset$. It is 
	vacuously true that $C_0$ is an order relation on $A_0$. Thus, we will assume that $A_0$ is non-empty.
	Let $x, y \in A_0$ such that $x \neq y$. Since $A_0 \subseteq A$ it must be that $x, y \in A$.
	Therefore, either $xCy$ or $yCx$ since $C$ is an order relation. 
	Note that $(x, y), (y, x) \in A_0 \times A_0$ and thus $C$ will contain either $(x, y)$ or 
	$(y, x)$ depending on whichever element $C$ contains. In either case, we can conclude that 
	$C_0$ is comparable in $A_0$. Now, let $x \in A_0$, which implies that $x \in A$. Since $C$ is an order
	relation on $A$, this implies that $(x, x) \notin C$. Hence, $(x, x) \notin C_0$ by construction.
	Therefore, $C_0$ must be non-reflexive. Let $(x, y), (s, t) \in C_0$. This implies 
	that $(x, y), (s, t) \in C$. Since $C$ is an order relation on $A$, $(x, t) \in C$ must be true.
	Since $x, t \in A_0$, $(x, t) \in A_0 \times A_0$ and therefore $(x, t) \in C_0$. Thus we conclude
	that $C_0$ is transitive and therefore an order relation.
\end{proof}

\exnumber{8.} Check that the relation defined in Example 7 is an order relation.

\begin{proof}
	The order relation in Example 7 is the following: For the set of real numbers, define $xCy$
	if $x^2 < y^2$, or if $x^2 = y^2$ and $x < y$.
	
	We begin by showing that this relation is comparable. Let $x, y \in \RR$ such that 
	$x \neq y$. Since Note either $x^2 < y^2$, $x^2 > y^2$ or $x^2 = y^2$. In the first
	two cases we can conclude that either $xCy$ or $yCx$ showing our relation is comparable. In
	the third case, we then check if $x < y$ or $x > y$ since $x \neq y$. Note that in either case
	we can conclude that $xCy$ or $yCx$ showing or relation is comparable. Hence, we conclude
	that this relation is comparable on the real numbers.
	
	Now we demonstrate that the relation is non-reflexive. Let $x \in \RR$, then $x^2 = x^2$ and 
	$x = x$. Note that based on our definition for our relation we cannot relate $x$ to itself. Thus,
	our relation does not compare elements to themselves and we can conclude that it is non-reflexive.
	
	To show that the relation is transitive, let $xCy$ and $yCz$. By definition,
	this implies that either $x^2 < y^2,$ or $x^2 = y^2$ and $x < y$. Similarly, we have that 
	$y^2 < z^2$, or $y^2 = z^2$ and $y < z$. If $x^2 < y^2$ and $y^2 < z^2$ note that $x^2 < z^2$
	must be true and we conclude that $xCz$. If $x^2 = y^2$ and $y^2 < z^2$, note that
	$x^2 < z^2$ must be true and we conclude that $xCz$. If $x^2 = y^2$ and $y^2 = z^2$, then note
	$x < y$ and $y < z$. This implies that $x < z$, so we conclude that $xCz$. The remaining case 
	is analogous to the one we did earlier but the signs are swapped. Hence, we conclude that
	our relation is transitive. Therefore, $C$ is an order relation on the real numbers.
\end{proof}

\exnumber{9.} Check that the dictionary order is an order relation. 

\begin{proof}
	Let $A$ and $B$ be sets such that $<_A$ and $<_B$ are order relations in their respective sets.
	We define the dictionary order on $A \times B$ as the relation where $(a, b) < (c, d)$
	if $a <_A c$, or if $a = c$ and $b <_B d.$ To show that the dictionary order is an
	order relation we must show that it has comparability, non-reflexivity, and transitivity.
	
	We begin by showing that this relation is comparable. Let $(a, b), (c, d) \in A \times B$
	such that $(a, b) \neq (c, d)$.
	If $a <_A c$, then $(a, b) < (c, d)$, and if $c <_A a$ then $(c, d) < (a, b)$. If 
	$a = c$, it must be that $b \neq c$ by construction. Thus, either $b <_B d$ or $d <_B b$.
	In either case we can compare both elements of $A \times B$. Therefore, we conclude that
	the dictionary order is an order relation.
	
	We now demonstrate that the dictionary order is non-reflexive. Let $(x, y) \in A \times B$.
	Then, $x = x$ and $y = y$. Note that since $y = y$, we cannot compare $(x, y)$ to itself
	by the definition of the dictionary order. Hence, an element in $A \times B$ is not comparable
	to itself under this relation. Therefore, we conclude that the dictionary order is 
	non-reflexive.
	
	Finally, to show that the dictionary order is transitive let $(a, b) < (c, d)$ and 
	$(c, d) < (n, m)$. By definition, either $a <_A c$ or $a = c$ and $b <_B d$. Similarly.
	either $c <_A n$ or $c = n$ and $d <_B m $. If $a <_A c$ and $c <_A n$, then 
	$a <_A n$ by the fact $<_A$ is an order relation in $A$. Thus we conclude that 
	$(a, b) < (n, m)$. If $a = c$ and $c = n$, it must be that $b <_B d$ and $d <_B m$. Since
	$<_B$ is an order relation in $B$ it must be that $b <_B m$. Thus, $(a, b) < (n, m)$. 
	If $a = c$ and $d <_B m$ it must be that $b <_B d$. This implies that $(a, b) < (n, m)$.
	If $a <_A c$ and $c = n$, we conclude that $a <_A n$, showing that $(a, b) < (n, m)$.
	Therefore, the dictionary order is transitive and we further conclude that it is an order 
	relation on $A \times B$.
\end{proof}

\exnumber{10.}
\begin{enumerate}[label=(\alph*)]
	\item Show that the map $f: (-1, 1) \to \RR$ of Example 9 is order preserving.
	
	\begin{proof}
		The function in Example 9 is the following: $f: (-1, 1) \to \RR$ such that 
		$$
		f(x) = \frac{x}{1 - x^2}.
		$$
		To show that $f$ is order preserving, we must demonstrate that for $x, y \in (-1, 1)$
		such that $x < y$, it must be that $f(x) < f(y)$. Thus, assume that $x, y \in (-1, 1)$
		such that $x < y$. We will break our analysis into cases, the first of which is 
		when $x \geq 0$ and $y > 0$. Consider the difference $y^2 - x^2$ which factors into
		$(y - x)(y + x)$. Since $y > x$, we have that $y - x > 0$. Given that $x \geq 0$ and $y > 0$
		we also have that $y + x > 0$. Therefore, $y^2 - x^2 > 0$ and we can conclude that 
		$y^2 > x^2$. This implies that $-y^2 < -x^2$, and furthermore $1 - y^2 < 1 - x^2$.
		Finally, note that
		$$
		\frac{1}{1 - y^2} > \frac{1}{1 - x^2}.
		$$
		Thus, we have the following statements:
		$$
		\frac{x}{1 - x^2} < \frac{y}{1 - x^2}
		$$
		since $x < y$, and 
		$$
		\frac{y}{1 - x^2} < \frac{y}{1 - y^2}
		$$
		due to our previous work. Therefore,
		$$
		\frac{x}{1 - x^2} < \frac{y}{1 - y^2},
		$$
		implying $f(x) < f(y)$ just as we want.
		
		Now, consider the case $x < 0$, $y \leq 0$. Then, note that $-x > 0$, $-y \geq 0$, and $-x > -y$.
		Such that $f(-x) > f(-y)$ by our previous work. Furthermore.
		$$
		f(-x) = \frac{-x}{1- x^2} = -f(x)
		$$
		$$
		f(-y) = \frac{-y}{1 - y^2} = -f(y).
		$$
		Thus, $-f(x) = f(-x) > f(-y) = -f(y)$ which implies $f(x) < f(y)$.
		
		Finally, now consider $x < 0$ and $y > 0$. Note that $f(x) < 0$ and $f(y) > 0$ by definition, and
		thus we conclude that $f(x) < f(y)$. Since we showed that $f(x) < f(y)$ when $x < y$ for 
		all of our cases we can conclude that $f$ is order preserving.
	\end{proof}
	 
	\item Show that the equation $g(y) = 2y/[1 + (1 + 4y^2)^{1/2}]$ defines a function
	$g: \RR \to (-1, 1)$ that is both a left and a right inverse for $f$.
	\begin{proof}
		We first show that $g$ is a left inverse, i.e., that $g(f(x)): (-1, 1) \to (-1, 1)$ is the identity
		map on $(-1, 1)$.
		We first have that 
		$$
		g(f(x)) = \frac{2f(x)}{1 + \big(1 + 4f(x)^2\big)^{1/2}}
		$$  
		by definition of $g$. Using the definition of $f$ we now have
		$$
		g(f(x)) = \frac{2x}{1 - x^2} \cdot \frac{1}{1 + \left(1 + 4\left(\frac{x}{1 - x^2}\right)^2\right)^{1/2}}
		$$
		$$
		g(f(x)) = \frac{2x}{1 - x^2} \cdot 
		\frac{1}{1 + \left(\frac{(1-x^2)^2 + 4x^2}{(1-x^2)^2}\right)^{1/2}}
		$$
		$$
		g(f(x)) = \frac{2x}{1 - x^2} \cdot \frac{1}{1 + \frac{x^2 - 1}{1 - x^2}}
		$$
		$$
		g(f(x)) = \frac{2x}{1 - x^2} \cdot \frac{1}{\frac{(1 - x^2) + x^2 + 1}{1 - x^2}} = 
		\frac{2x}{1 - x^2} \cdot \frac{1 - x^2}{2}
		$$
		$$
		g(f(x)) = x
		$$
		Thus we conclude that $g$ is a right inverse for $f$.
		
		Now, we will check that $g$ is a right inverse of $f$, i.e., that $f(g(y)): \RR \to \RR$ is the identity
		map on $\RR$.
		
		We first have that 
		$$
		f(g(y)) = \frac{g(y)}{1 - g(y)^2}
		$$
		by the definition of $f$, and using the definition of $g$ we get
		$$
		f(g(y)) = \frac{\frac{2y} {1 + (1 + 4y^2)^{1/2}}}{1 - \left(\frac{2y}{1 + (1 + 4y^2)^{1/2}}\right)^2}
		$$
		$$
		f(g(y)) = \frac{2y}{1 + (1 + 4y^2)^{1/2}} \cdot \frac{1}{1 - \frac{4y^2}{4y^2 + 2 + 
		2(1 + 4y^2)^{1/2}}}
		$$
		$$
		f(g(y)) = g(y) \cdot 
		\frac{1} {\frac{4y^2 + 2 + 2(1 + 4y^2)^{1/2} - 4y^2} {4y^2 + 2 + 2(1 + 4y^2)^{1/2}}}
		$$
		$$
		f(g(y)) = g(y) \cdot 
		\frac{4y^2 + 2 + 2(1 + 4y^2)^{1/2}}{2 + 2(1 + 4y^2)^{1/2}}
		$$
		$$
		f(g(y)) = g(y) \cdot 
		\frac{2y^2 + 1 + (1 + 4y^2)^{1/2}}{1 + (1 + 4y^2)^{1/2}}
		$$
		$$
		f(g(y)) = g(y) \cdot \left(\frac{2y^2}{1 + (1 + 4y^2)^{1/2}} + 1 \right)
		$$
		$$
		f(g(y)) = g(y) \cdot \left(yg(y) + 1\right) = yg(y)^2 + g(y)
		$$
		$$
		f(g(y)) = \frac{4y^3}{1 + 2(1 + 4y^2)^{1/2} + (1 + 4y^2)} + \frac{2y}{1 + (1 + 4y^2)^{1/2}}
		$$
		$$
		f(g(y)) = \frac{4y^3 + 2y(1 + (1 + 4y^2)^{1/2})}{1 + 2(1 + 4x^2)^{1/2} + (1 + 4y^2)}
		$$
		$$
		f(g(y)) = \frac{2\left(2y^3 + y + y(1 + 4y^2)^{1/2}\right)}{2 \left(1 + (1 + 4y^2)^{1/2} + 2y^2\right)}
		$$
		$$
		f(g(y)) = \frac{y\left(2y^2 + (1 + (1 + 4y^2)^{1/2})\right)}{2y^2 + (1 + (1 + 4y^2)^{1/2})}
		$$
		$$
		f(g(y)) = y
		$$
		Thus, we conclude that $g$ is also a right inverse of $f$.
	\end{proof}
\end{enumerate}

\exnumber{11.} Show that an element in an ordered set has at most one immediate successor 
and at most one immediate predecessor. Show that a subset of an ordered set has at most
one smallest element and at most one largest element.

\begin{proof}
	Let $S$ be a non-empty ordered set and $a \in S$ be an element of $S$. If $b \in S$ is an immediate
	successor of $a$, then by definition the set $B = \{x \in S \mid a < x < b\}$ is empty. Assume that
	$c$ is another immediate successor of $a$. Then, similarly, the set $C = \{x \in S \mid a < x < c\}$
	is empty. Since $S$ is an ordered set, if $b \neq c$, then either $b < c$ or $c < b$. Note that 
	if $b < c$, then $b \in C$ which is a contradiction of $c$ being an immediate successor of $a$.
	Similarly, if $c < b$ it implies that $c \in B$ which is a contradiction of $b$ being an immediate
	successor of $a$. Therefore, it must be that $b = c$. As such, $a$ can have at most one
	immediate successor. 
	
	As with the case of the immediate successor of $a$, we can show in an analogous argument that 
	$a$ can have at most one immediate predecessor. As such, let $b$ and $c$ be immediate predecessors
	of $a$. By definition, we have that the sets $B = \{x \in S \mid b < x < a\}$ and 
	$\{C = x \in S \mid c < x < a\}$ are empty. Given $S$ is an ordered set, if $b \neq c$ it must
	be that either $b < c$ or $c < b$. If $b < c$ then $c \in B$ which is a contradiction to $c$
	being an immediate predecessor of $a$. If $c < b$ then $b \in C$ which is a contradiction to
	$b$ being an immediate predecessor of $a$. Therefore, it must be that $b = c$, showing that 
	$a$ can have at most one immediate predecessor. 
	
	Let $A \subseteq S$ be a subset of $S$. If $A$ is empty, then it has no smallest or 
	largest element, given it has no elements at all. By \exnumber{7} we have that 
	the restriction of the order relation on $S$ in respect to $A$ is also an order relation.
	Let $a \in A$ be a largest element of $S$, this implies that for any $c \in A$ we must 
	have that $c < a$. If $b \in A$ were another largest element of $A$ we would have that 
	$b < a$ and $a < b$ are both true, which is a contradiction to $A$ being an ordered set.
	Therefore, if both $a$ and $b$ are the largest element of $A$ it must be that $a = b$.
	Thus, $A$ can have at most one largest element.
	
	Similar with the case of a largest element, we can show that $A$ can have at most one smallest 
	element. Let $a \in A$ be a smallest element, then for any $c \in A$ it must be true 
	that $a < c$. If $b \in A$ were another smallest element of $A$, then both $a < b$ and 
	$b < a$ are true, which contradicts $A$ being an ordered set. Therefore, $a = b$ which implies
	that $A$ can have at most one smallest element. 
\end{proof}

\exnumber{12.} Let $\ZZ_+$ denote the set of positive integers. Consider the following 
order relations on $\ZZ_+ \times \ZZ_+$:
\begin{enumerate}[label=(\roman*)]
	\item The dictionary order.
	\item $(x_0, y_0) < (x_1, y_1)$ if either $x_0 - y_0 < x_1 - y_1$, or 
	$x_0 - y_0 = x_1 - y_1$ and $y_0 < y_1$.
	\item $(x_0, y_0) < (x_1, y_1)$ if either $x_0 + y_0 < x_1 - y_1$, or 
	$x_0 + y_0 = x_1 + y_1$ and $y_0 < y_1$.
\end{enumerate}
In these order relations, which elements have immediate predecessors? Does the set have 
a smallest element? Show that all three order types are different. 

As an addendum I noticed late into this that $\ZZ_+$ does not include $0$. However, this 
doesn't change much other than what are the smallest elements, e.g. $(1,1)$ instead of $(0,0)$,
or that the elements with no immediate predecessor are actually of the form $(1, y)$ instead
of $(0, y)$. However, the analysis of each order relation on the set is the same, with the core
facts remaining intact so I rather not rewrite this and risk causing an error along the way. So, 
for the purpose of this exercise simply assume I'm talking about the set
$\ZZ_{\geq 0} \times \ZZ_{\geq 0}$.

We first consider the dictionary order. Under this order, note that there is a smallest
element in $\ZZ_+ \times \ZZ_+$, which is $(0,0)$. This is due to the fact that for
any $(x, y)$ if $x \neq 0$ then $0 < x$, so $(0,0) < (x, y)$. If $x = 0$, then $y \neq 0$ which implies
that $0 < y$ making $(0, 0) < (x, y)$. Therefore, $(0,0) < (x, y)$ for all $(x, y)$ in this set.
When it comes to which elements have an immediate predecessor and which don't under this order
relation, consider the following:

Elements of the form $(x, 0)$ cannot have an immediate predecessor. To show this, consider the kinds
of elements that are smaller than $(x, 0)$. If $(y,z)$ is smaller that $(x, 0)$ either 
$y < x$ or $y = x$ and $z < 0$. Note that $z$ cannot be smaller than $0$ which implies
that $(y, z)$ can only be smaller than $(x, 0)$ if $y < x$. In particular, any point of the form
$(x - 1, z)$ will be smaller than $(x, 0)$. As such, if $(x - 1, z)$ were the immediate predecessor 
of $(x, 0)$ then the set $\{(a, b) \mid (x - 1, z) < (a, b) < (x, 0)\}$ is empty. But, note that 
$(x - 1, z) < (x - 1, z + 1)$ and $(x - 1, z + 1) < (x, 0)$. Since this is true for any $z$ we 
conclude that $(x, 0)$ cannot have an immediate predecessor under the dictionary order.
Any other element of this set (except the smallest element $(0,0)$ also) does have an immediate 
predecessor. In particular for the element $(x, y)$ where $y \neq 0$, the element $(x, y - 1)$
is its immediate predecessor. 

Now, we consider the order relation in (ii). We first note that there can be no smallest
element with this order. To show this, assume that $(x, y)$ was the smallest element 
under this order. Thus, for any $(a, b)$ we have that $(x, y) < (a, b)$. If $x \neq 0$,
note that the element $(x - 1, y) < (x, y)$ by how we defined our order. Thus, it must be 
that $x = 0$. However, note that for any $y$, we have $y + 1$ such that $(0, y) < (0, y+1)$ by 
definition of this order relation. As such, it cannot be that $(x, y)$ is the smallest element.

In regard to immediate predecessors, we note that under this order relation, elements of the 
form $(x, 0)$ and $(0, y)$ do not have an immediate predecessor. To see this, consider 
$(x, 0)$ where $x \neq 0$, and assume that $(a, b)$ is its immediate predecessor. Then, either
$a - b < x$ or $a - b = x$ and $b < 0$. Since $b$ cannot be less than $0$ it must 
be that $a - b < x$. In particular, it must be that $a - b = x - 1$ as it were smaller it could not
be the immediate predecessor of $(x, 0)$. Note then, that any element of the form $(a + n, b + n)$
where $n$ is any positive integer will have both a difference of $x - 1$ and a component such that 
$b < b + n$. As such, $(a, b) < (a + n, b + n) < (x, 0)$ which is a contradiction to
$(a, b)$ being the immediate predecessor of $(x, 0)$. For elements of the form $(0, y)$ we can 
construct an analogous argument as to why they also do not have an immediate predecessor, including 
the element $(0, 0)$. However, every other element does have an immediate predecessor. In particular
for an element $(x, y)$ where $x, y \neq 0$ its immediate predecessor is $(x - 1, y - 1)$.
 
For (iii) we note that under this order relation there is a smallest element, namely $(0,0)$, 
and every other element will have an immediate predecessor. 
To show this, simply note that $0 + 0 = 0$, and that for an element $(x, y)$ if either
$x$ or $y$ is non-zero then their sum will be greater than $0$, as such $(0,0)$ must be 
the smallest element. In regard to immediate predecessors, consider the element $(x, y)$ such
that it is not $(0,0)$. If $y = 0$ then  note $x + y = x$. Since $y = 0$, it must be that $(x, y)$
is the smallest element whose component sum equals $x$, thus the element $(0, x - 1)$ will be its 
immediate predecessor, given it is the greatest element whose sum is one less than $x$.
For the case where $y \neq 0$, note that the immediate predecessor of $(x, y)$ will be 
the element $(x + 1, y - 1)$. 

From our analysis, we note that under the dictionary order there is a smallest element $(0,0)$
but elements of the form $(x, 0)$ will have no immediate predecessor. Under the order relation
in (ii) there will be no smallest element and elements of the form $(x, 0)$ and $(0, y)$ will
have no immediate predecessor. Finally, under the order relation of (iii) there will be a smallest element
$(0,0)$ and every other element will have an immediate predecessor. As such, all of these order
relationships must be different relations which can be easily seen by the different kinds of 
ordered pairs that are not included in each relation due to their difference in which elements
have or don't have an immediate predecessor. 

\exnumber{13.} Prove the following:
\begin{theorem}
	If an ordered set $A$ has the least upper bound property, then it has the 
	greatest lower bound property.
\end{theorem}

\begin{proof}
	Let $A$ be an ordered set with the least upper bound property. Consider the subset
	$B_0 \subseteq A$ which is bounded below. Since $B_0$ is bounded below there exists
	some $a \in A$ such that for all $b \in B_0$ we have that $a \leq b$. Thus, consider 
	the set of all lower bounds of $B$, and call it $A_0$. Note that $A_0$ is bounded
	above, since for all $a \in A$ we have that $a \leq b$ for some $b \in B_0$.
	In particular, given that $A$ has the least upper bound property, it must be 
	that $A_0$ has a least upper bound, say $a_0$. If $a_0 \notin A_0$, then 
	there exists some $b \in B_0$ such that $b < a_0$, but note that $b$ is an upper
	bound of $A_0$ which contradicts the fact that $a_0$ is the least upper bound of 
	$A_0$. Therefore, $a_0$ must be an element of $A_0$. In particular, this implies
	that $a_0$ is a lower bound of $B_0$. Furthermore, since $a_0$ is the largest element of 
	$A_0$ it is also the greatest lower bound of $B_0$. Since $B_0$ was an arbitrary 
	set in $A$ that was bounded below we conclude that $A$ has the greatest lower bound property
	also.
\end{proof}

\exnumber{14.} If $C$ is a relation on a set $A$, define a new relation $D$ on $A$ 
by letting $(b, a) \in D$ if $(a, b) \in C$.
\begin{enumerate}[label=(\alph*)]
	\item Show that $C$ is symmetric if and only if $C = D$.
	\begin{proof}
		Let $C$ be symmetric, then if $(x, y) \in C$ we also have that $(y, x) \in C$.
		Given we defined $D$ such that $(y, x) \in D$ if $(x, y) \in C$ we note that 
		$C = D$. This is because since $(x, y), (y, x) \in C$, we also have that 
		$(x,y), (y, x) \in D$.
		
		Now, let $C = D$.
		Let $(x, y) \in C$, by our definition of $D$ we have that $(y, x) \in D$.
		Since $C = D$, we have that $(y, x) \in C$. Since $(x, y) \in C$ was arbitrary,
		we note that $C$ is symmetric.
	\end{proof}
	\item Show that if $C$ is an order relation, $D$ is also an order relation.
	\begin{proof}
		Let $C$ be an order relation. This means that $C$ is comparable, non-reflexive, and transitive.
		Since $C$ is comparable, we have that for any $x, y \in A$ such that $x \neq y$, either
		$xCy$ or $yCx$. In either case, due to how we defined $D$ we have $yDx$ or $xDy$ respectively.
		As such, we conclude that $D$ is comparable. Since $C$ is non-reflexive we have that 
		for any $x \in A$, the relation $xCx$ does not hold. Similarly, this implies that 
		$xDx$ also does not hold. As such, $D$ is non-reflexive. Given that $C$ is transitive,
		if $xCy$ and $yCz$, then $xCy$. Based on our definition of $D$, if $xCy$ and $yCz$
		we have $yDx$ and $zDy$, and similarly if $xCz$, then $zDx$. Thus, $D$ is also transitive.
		Therefore, $D$ is also an order relation.
	\end{proof}
	\item Prove the converse of the theorem in Exercise 13.
	\begin{proof}
		The converse of the theorem in Exercise 13, would be:
		
		\begin{theorem}
			If an ordered set $A$ has the greatest lower bound property, then it has the least
			upper bound property.
		\end{theorem}
		
		The proof will be similar to that of Exercise \exnumber{13}. Let $(S,<)$ be an ordered set
		with the greatest lower bound property. Consider $A \subseteq S$ that us bounded above.
		Since $A$ is bounded above, there exists $b \in S$ such that $a \leq b$ for all $a \in A$.
		Let $B$ be the set of all upper bounds of $A$, and note it is not empty since $b$ exists.
		In particular, $B$ is bounded below since for any $a \in A$ we have that $a \leq b$.
		Given that $S$ has the greatest lower bound property, $B$ has a greatest lower bound,
		namely $b_0$. If $b_0$ is not an element of $B$, then $b_0$ is not an upper bound of $A$. 
		Thus, there exists some $a \in A$ such that $b_0 < a$. This, however, is 
		a contradiction since any $b \in B$ is defined such that $a \leq b$ for all $a \in A$. 
		Therefore, $b_0 \in B$ and we conclude that $b_0$ is an upper bound of $A$. Furthermore, given
		$b_0$ is the greatest lower bound of $B$, we have that $b_0 \leq b$ for all $b \in B$ and thus
		is also the least upper bound of $A$. Since $A$ was any subset of $S$ that is bounded
		above we conclude that $S$ has the least upper bound property.
	\end{proof}
\end{enumerate}

\exnumber{15.} Assume that the real line has the least upper bound property.
\begin{enumerate}[label=(\alph*)]
	\item Show that the sets
	$$
	[0, 1] = \{x \mid 0 \leq x \leq 1\},
	$$
	$$
	[0, 1) = \{x \mid 0 \leq x < 1\}
	$$
	have the least upper bound property. 
	\begin{proof}
		Consider the subsets $(a, b) \subseteq [0,1]$. Note that $b$ is the least upper bound
		of $(a, b)$, which would be the same if the subset was $(a,b]$. For subsets of the 
		form $[a, b]$ and $[a, b)$ we are still left with $b$ being the least upper bound.
		Since every subset that is bounded above has a least upper bound we conclude the set 
		$[0, 1]$ has the least upper bound property. 
		
		Thus, consider subsets of $[0, 1)$. In particular, subsets of the form $[a, b]$ or
		$(a, b]$ note 
		that $b$ cannot equal $1$ as that is not included as an element of the set. Given this 
		restriction on $b$ we note that $b$ will remain the least upper bound of the subset. 
		For subsets of the form $(a, b)$ and $[a, b)$ we first note that if $b = 1$ then the 
		subset is not bounded above, as such those do not have a least upper bound given they have
		no upper bound at all. If $b \neq 1$ then one simply notices that $b$ will be the least 
		upper bound of the subset. Thus, since every subset that is bounded above has a least
		upper bound the set $[0, 1)$ has the least upper bound property. 
 	\end{proof}
	\item Does $[0,1] \times [0, 1]$ in the dictionary order have the least upper
	bound property? What about $[0,1] \times [0,1)$?
	What about $[0,1) \times [0, 1]$?
	
	We begin by first showing that $[0, 1] \times [0, 1)$ does not have the least upper bound
	property. Consider the non-empty subset $\{0\} \times [0, 1)$. Note that is is bounded above
	since the point $(1, 0)$ is greater than any point on the subset by our dictionary order relation.
	As such, $\{0\} \times [0, 1)$ is a non-empty subset of our set that is also bounded above. Assume for 
	contradiction that the point $(x_0, x_1)$ is the least upper bound of the subset. If $x_0 = 0$, then
	it must be that $x_1 \geq b$ for $b \in [0, 1)$ such that $(0, b) \in \{0\} \times [0, 1)$. Note, however,
	that this cannot be, since $x_1 \in [0, 1)$ and thus we can define $b$ as $\frac{1 + x}{2}$ and see that
	$b < 1$ and $b > x$, as such $(0, b) \in \{0\} \times [0, 1)$ and $(0, x_1) < (0, b)$ which is a contradiction.
	Thus, it must be that $x_0 \neq 0$. However, if $x_0 = a$ for $a > 0$, we can define $c = \frac{a}{2} > 0$
	which is smaller than $a$, and have that $(c, 0)$ is an upper bound on our subset, and also that 
	$(c, 0) < (a, x_1)$ which is a contradiction to the fact $(x_0, x_1)$ is the least upper bound 
	on out subset. Therefore, it cannot be that $(x_0, x_1)$ is a least upper bound, and thus
	our subset has no least upper bound, preventing this ordered set from having the least upper bound
	property. 
	
	We claim that the ordered set $[0, 1] \times [0,1]$ does have the least upper bound property under 
	the dictionary order relation. To show this, consider a subset $S \subseteq [0, 1] \times [0,1]$ that 
	is non-empty and bounded above. Then, there must exist a point $(x_0, x_1) \in [0, 1] \times [0,1]$
	such that $(x_0, x_1) \geq (a, b)$ for any $(a, b) \in S$. Let $S_x$ be defined as
	$\{x \in [0, 1] \mid (x, y) \in S, y \in [0, 1]\}$, i.e., $S_x$ is the set of all $x \in [0, 1]$
	such that $(x, y) \in S$ for at least some $y \in [0, 1]$. Then, note that since $(x_0, x_1)$ is an 
	upper bound on $S$, that $x_1$ is also an upper bound on $S_x$, since by definition $x_1 \geq x$ for 
	all $x \in S_x$. Thus, $S_x$ is a non-empty bounded above subset of $[0,1]$ and therefore must have 
	a least upper bound on $[0, 1]$ by part (a). Let $\alpha$ be such least upper bound. If $\alpha \notin S_x$,
	then note $(\alpha, 0)$ is the least upper bound of $S$. We know that $(\alpha, 0)$ is an upper bound of 
	$S$ since $\alpha > x$ for all $x \in S_x$. Furthermore, it is also a least upper bound, since for any 
	element less than $(\alpha, 0)$, it must be that their first component, say $x_0$ is less than $\alpha$.
	So, if $(x_0, x_1)$ is an upper bound of $S$ such that $x_0 < \alpha$ we land on a contradiction, since 
	$\alpha$ is the least upper bound of $S_x$, it must be that there exists at least some $(x, y) \in S$
	such that $x > x_0$. Thus, $(\alpha, 0)$ is the least upper bound of $S$.
	
	If $\alpha \in S_x$, then consider the analogous set $S_y$ defined as the set of all $y \in [0, 1]$
	such that $(\alpha, y) \in S$ for at least some $x \in [0, 1]$. Since $\alpha \in S_x$, we note that 
	there must exists at least one $(\alpha, y) \in S$, and thus we conclude that $S_y$ is non-empty.
	Furthermore, given that $S_y$ is a subset of $[0, 1]$ we note that it is bounded above by $1$, therefore,
	$S_y$ is a non-empty bounded above subset of $[0, 1]$. Thus, $S_y$ has a least upper bound in $[0, 1]$
	by part (a), call it $\beta$. We claim that $(\alpha, \beta)$ is the least upper bound of $S$.
	Firstly, we can see that $(\alpha, \beta)$ is an upper bound of $S$ by our construction of the point, thus
	we only need to examine if it is the least upper bound. To see it is the least upper bound, consider
	the point $(c, d)$ and say it is an upper bound of $S$ such that $(c, d) < (\alpha, \beta)$. Note
	it cannot be that $c < \alpha$, since $\alpha$ is the least upper bound of $S_x$ which implies that 
	any point smaller that $\alpha$, including $c$, must be smaller that at least one $x \in S_x$, implying that 
	$(c, d) < (x, y)$ for some $(x, y) \in S$. Thus, $c = \alpha$. Then, it must be that 
	$d < \beta$ which implies that $\beta \neq 0$ and thus that $\alpha \in S_x$ by our previous work. In particular,
	since $\beta$ is the least upper bound of $S_y$, it must be that $d$ is not an upper bound of $S_y$.
	Thus, there exists some $(\alpha, y) \in S$ such that $(\alpha, y) > (\alpha, d)$. Hence, we conclude
	that $(c, d)$ cannot be an upper bound of $S$ if it is smaller than $(\alpha, \beta)$. Thus, we 
	have that $(\alpha, \beta)$ is the least upper bound of $S$. Given that $S$ was an arbitrary non-empty
	and bounded above subset of $[0, 1] \times [0, 1]$ we conclude that $[0, 1] \times [0,1]$ has the
	least upper bound property.
	
	In regard to $[0,1) \times [0,1]$ we also claim that it has the least upper bound property. The proof
	of this claim is analogous to the one we just made for $[0, 1] \times [0,1]$. In fact, the logic
	of our proof for that set is the exact same one used for this set. The only reason this logic fails
	for $[0,1] \times [0,1)$ is because the set $[0,1)$ for the second component prevents us from 
	asserting definitively that a subset $S_y$ is bounded above, as in the case of our counter-example. 
	\end{enumerate} 
 
%%%%%%%%%%%%%%%%%%%%%%%%%%%%%%%%%%%%%%%%%%%%%%%%%%%%%%%%%%%%%%%%%%%%%%
\end{document}%%%%%%%%%%%%%%%%%%%%%%%%%%%%%%%%%%%%%%%%%%%%%%%%%%%%%%%%
%%%%%%%%%%%%%%%%%%%%%%%%%%%%%%%%%%%%%%%%%%%%%%%%%%%%%%%%%%%%%%%%%%%%%%
